\documentclass[10pt,a4paper]{amsart}
\usepackage[margin=19mm]{geometry}
\usepackage{amsmath,amssymb,mathtools}
\usepackage[scr=rsfs,cal=euler]{mathalfa}
\usepackage{xfrac}
\usepackage[unicode,hidelinks,bookmarks=false]{hyperref}
\newcommand{\ie}{i.e.\ }
\newcommand{\cf}{cf.\ }
\newcommand{\resp}{respectively\ }
\newcommand{\Subsection}{\subsection}
\providecommand{\nobreakdash}{\nobreak\hbox{-}}
\setlength{\emergencystretch}{2em}
\multlinegap0pt
\usepackage{lmodern}
\usepackage{mathtools}
\usepackage{mathrsfs}
\usepackage{bm}
\usepackage{tikz-cd}
\usepackage{enumitem}
\usepackage{doi}
\newtheorem{theorem}{Theorem}[section]
\newtheorem{proposition}[theorem]{Proposition}
\newtheorem{lemma}[theorem]{Lemma}
\newtheorem{corollary}[theorem]{Corollary}
\newtheorem{definition}[theorem]{Definition}
\newtheorem{assumption}[theorem]{Assumption}
\newtheorem{example}[theorem]{Example}
\newtheorem{remark}[theorem]{Remark}
\newcommand{\Hilb}{\mathcal H}
\newcommand{\Dom}{\mathcal D}
\newcommand{\Ran}{\operatorname{Ran}}
\newcommand{\Tr}{\operatorname{Tr}}
\newcommand{\FP}{\operatorname{FP}}
\newcommand{\Id}{\mathbf 1}
\newcommand{\ii}{\mathrm i}
\newcommand{\ee}{\mathrm e}
\newcommand{\dd}{\,\mathrm d}
\newcommand{\R}{\mathbb R}
\newcommand{\N}{\mathbb N}
\newcommand{\C}{\mathbb C}
\newcommand{\Heff}{H_{\mathrm{eff}}}
\newcommand{\HFM}{H_{\mathrm{FM}}}
\newcommand{\PN}{P_N}
\newcommand{\UN}{U_N}
\newcommand{\CinftyH}{C^\infty(H_0)}
\newcommand{\Op}{\operatorname{Op}}
\begin{document}
\title[Spectral Cut-off Oscillatory Integrals for Non-Autonomous Ha]{Spectral Cut-off Oscillatory Integrals for Non-Autonomous Hamiltonian Evolution Equations}
\author{Jean-Pierre Magnot $^*$ Univ. Angers, CNRS, LAREMA, SFR MATHSTIC, F-49000 Angers, France; and; Lycée Jeanne d'Arc,; Avenue de Grande Bretagne,; Secondary 81Q05, 81S40, 47A58, 58J40 ęywords Non-autonomous Schrödinger equations; unbounded Hamiltonians; oscillatory integrals; Jean-Pierre Magnot}
\date{}
\maketitle
\noindent\emph{Source and licence.} Spectral Cut-off Oscillatory Integrals for Non-Autonomous Hamiltonian Evolution Equations. Version arXiv:2605.26899v2 (CC BY 4.0). This material is distributed under the Creative Commons Attribution 4.0 International licence, http://creativecommons.org/licenses/by/4.0/, as recorded in the arXiv OAI licence metadata for this exact version and shown on the abstract page https://arxiv.org/abs/2605.26899v2. Full rights notice: licenses/arxiv-2605.26899-CC-BY-4.0.txt.\par\smallskip
\noindent\textit{Editorial scope.} Counted unit: the original \texttt{theorem} environment \texttt{thm:construction-by-spectral-cutoff} at source lines 1897--1955 together with its complete proof at lines 1957--2132; the delivered document keeps source lines 407--2150 so that every referenced label and every citation resolves inside it. Native editable source is the paper's own concatenated TeX; the AMS wrapper is editorial formatting only. No publisher class, upstream script or unrelated artwork is redistributed.
\label{sec:cutoff-oscillatory-integrals}

The aim of this section is to introduce the finite-dimensional
time-sliced objects from which the solutions will be constructed.  The
basic idea is simple: instead of postulating an infinite-dimensional
Feynman integral, we first project the Hamiltonian evolution onto
finite-dimensional spectral subspaces of a reference operator.  On
each such subspace, the propagator admits an exact finite-dimensional
state-sum representation.  When an additional configuration-space,
phase-space, or coherent-state realization is available, the same
time-sliced operators may also admit oscillatory integral
representations.  The full infinite-dimensional solution will then be
recovered by removing the spectral and temporal cut-offs.

\subsection{Reference operator and spectral cut-offs}

Let \(\Hilb\) be a complex separable Hilbert space.  We fix a positive
self-adjoint operator
\[
        H_0:D(H_0)\subset \Hilb\longrightarrow \Hilb.
\]
Throughout this section we assume that \(H_0\) has compact resolvent.
Thus its spectrum consists of a non-decreasing sequence of eigenvalues,
repeated according to multiplicity,
\[
        0\leq \lambda_1\leq \lambda_2\leq \cdots,
        \qquad
        \lambda_j\to+\infty,
\]
and \(\Hilb\) admits an orthonormal basis of eigenvectors of \(H_0\).

For \(N>0\), we denote by
\[
        P_N=\mathbf 1_{[0,N]}(H_0)
\]
the spectral projection of \(H_0\) associated with the interval
\([0,N]\).  We set
\[
        \Hilb_N=P_N\Hilb.
\]
Since \(H_0\) has compact resolvent, \(\Hilb_N\) is finite-dimensional.
The projections \(P_N\) satisfy
\[
        P_Nu\longrightarrow u
\]
strongly in \(\Hilb\) as \(N\to+\infty\).

The reference operator \(H_0\) also defines a scale of Hilbert spaces.
For \(s\geq0\), we set
\[
        \Hilb^s=D\big((1+H_0)^s\big),
        \qquad
        \|u\|_s=\|(1+H_0)^s u\|.
\]
We also write
\[
        \Hilb^\infty
        =
        C^\infty(H_0)
        =
        \bigcap_{s\geq0}\Hilb^s.
\]
The space \(\Hilb^\infty\) will play the role of a common smooth core.

\begin{remark}
The compact-resolvent assumption is natural for the construction of
finite-dimensional spectral cut-offs.  It can be relaxed if one replaces
the projections \(P_N\) by suitable finite-rank approximations or by
localized spectral cut-offs.  We keep the compact-resolvent setting in
order to isolate the main mechanism.
\end{remark}

\subsection{Cut-off Hamiltonians}

Let \(I\subset\mathbb R\) be a time interval, and let
\[
        t\in I\longmapsto H(t)
\]
be a family of symmetric operators on \(\Hilb\).  We assume, at this
stage, only that
\[
        \Hilb^\infty\subset D(H(t)),
        \qquad
        H(t)\Hilb^\infty\subset \Hilb
\]
for every \(t\in I\), and that the maps
\[
        t\longmapsto H(t)u
\]
are continuous for every \(u\in \Hilb^\infty\).  Stronger assumptions
will be introduced later in order to prove convergence of the cut-off
dynamics.

For each \(N\), we define the cut-off Hamiltonian
\[
        H_N(t)=P_NH(t)P_N
\]
as an operator on \(\Hilb_N\).  Since \(\Hilb_N\) is finite-dimensional,
\(H_N(t)\) is a bounded operator on \(\Hilb_N\).  If \(H(t)\) is
symmetric, then \(H_N(t)\) is self-adjoint on \(\Hilb_N\).  Therefore
the finite-dimensional non-autonomous Schr\"odinger equation
\[
        \ii\partial_t u_N(t)=H_N(t)u_N(t),
        \qquad
        u_N(s)=u_{N,s}\in \Hilb_N,
\]
has a unique unitary propagator
\[
        U_N(t,s):\Hilb_N\to\Hilb_N.
\]
It is characterized by
\[
        U_N(s,s)=\Id_{\Hilb_N},
        \qquad
        U_N(t,r)U_N(r,s)=U_N(t,s),
\]
and
\[
        \ii\partial_t U_N(t,s)u=H_N(t)U_N(t,s)u.
\]

For an initial datum \(u_0\in\Hilb\), the cut-off solution is
\[
        u_N(t)=U_N(t,0)P_Nu_0.
\]
The main convergence problem, treated in the next sections, is to find
conditions under which
\[
        u_N(t)\longrightarrow u(t)
\]
strongly in \(\Hilb\), uniformly for \(t\) in compact intervals, where
\(u(t)\) solves the original Hamiltonian equation
\[
        \ii\partial_t u(t)=H(t)u(t),
        \qquad
        u(0)=u_0.
\]

\subsection{Time-sliced products}

Fix \(N\) and let
\[
        s=t_0<t_1<\cdots<t_M=t
\]
be a partition of the interval \([s,t]\).  We write
\[
        \Delta t_j=t_{j+1}-t_j.
\]
The time-sliced approximation of \(U_N(t,s)\) is defined by
\[
        U_{N,M}(t,s)
        =
        \exp\big(-\ii\Delta t_{M-1}H_N(t_{M-1})\big)
        \cdots
        \exp\big(-\ii\Delta t_0H_N(t_0)\big).
\]
Equivalently,
\[
        U_{N,M}(t,s)
        =
        \prod_{j=M-1}^{0}
        \exp\big(-\ii\Delta t_jH_N(t_j)\big),
\]
where the product is ordered from right to left in increasing time.

Since \(H_N(t)\) is a continuous family of matrices, the standard
finite-dimensional theory of non-autonomous linear systems gives
\[
        U_{N,M}(t,s)\longrightarrow U_N(t,s)
\]
as the mesh of the partition tends to zero:
\[
        |\Pi_M|
        :=
        \max_{0\leq j\leq M-1}\Delta t_j
        \longrightarrow0.
\]
Thus, for fixed \(N\), the propagator \(U_N(t,s)\) is already obtained
as the limit of elementary oscillatory factors.

\subsection{Finite-dimensional state-sum representation}
\label{subsec:finite-dimensional-state-sum}

The time-sliced product introduced above admits a canonical finite-dimensional
representation as a sum over intermediate states. This representation should
be distinguished from a configuration-space oscillatory integral, which
requires additional geometric structure.

Let
\[
d_{N}=\dim\mathcal{H}_{N},
\]
and choose an orthonormal basis
\[
\mathcal{B}_{N}=(e_{N,1},\ldots,e_{N,d_{N}})
\]
of \(\mathcal{H}_{N}\). For a partition
\[
\Pi_{M}=\{s=t_{0}<t_{1}<\cdots<t_{M}=t\},
\]
set
\[
E_{N,j}
=
\exp\bigl(-\mathrm{i}\Delta t_{j}H_{N}(t_{j})\bigr),
\qquad
0\leq j\leq M-1.
\]
Thus
\[
U_{N,M}(t,s)
=
E_{N,M-1}\cdots E_{N,0}.
\]

For \(a,b\in\{1,\ldots,d_{N}\}\), write
\[
E_{N,j}(a,b)
=
\langle e_{N,a},E_{N,j}e_{N,b}\rangle.
\]
Repeated insertion of the identity
\[
I_{\mathcal{H}_{N}}
=
\sum_{a=1}^{d_{N}}
|e_{N,a}\rangle\langle e_{N,a}|
\]
gives the following exact formula:
\[
\begin{split}
&\langle e_{N,a_{M}},U_{N,M}(t,s)e_{N,a_{0}}\rangle
\\
&\quad=
\sum_{a_{1},\ldots,a_{M-1}=1}^{d_{N}}
\prod_{j=0}^{M-1}
E_{N,j}(a_{j+1},a_{j}).
\end{split}
\]

\begin{definition}[Spectral cut-off state-sum amplitude]
\label{def:spectral-cutoff-state-sum}
For fixed \(N\), \(M\), and endpoints \(a_{0},a_{M}\), the quantity
\[
\mathcal{A}_{N,M}(a_{M},a_{0};t,s)
=
\sum_{a_{1},\ldots,a_{M-1}=1}^{d_{N}}
\prod_{j=0}^{M-1}
E_{N,j}(a_{j+1},a_{j})
\]
is called the spectral cut-off state-sum amplitude associated with the basis
\(\mathcal{B}_{N}\) and the partition \(\Pi_{M}\).
\end{definition}

By construction,
\[
\mathcal{A}_{N,M}(a_{M},a_{0};t,s)
=
\langle e_{N,a_{M}},U_{N,M}(t,s)e_{N,a_{0}}\rangle.
\]
The state-sum is therefore an exact representation of the time-sliced
propagator and not an additional analytic object.

\begin{remark}
\label{rem:basis-dependence-state-sum}
The intermediate state-sum depends on the choice of the orthonormal basis
\(\mathcal{B}_{N}\), whereas the resulting operator \(U_{N,M}(t,s)\) does not.
In an abstract finite-dimensional Hilbert space, insertion of an orthonormal
basis produces a finite sum, not an integral over
\(\mathbb{R}^{d_{N}}\). A continuous configuration-space integral requires a
realization of \(\mathcal{H}_{N}\) as a space of functions, or a continuous
resolution of the identity such as a coherent-state representation.
\end{remark}

For \(u\in\mathcal{H}\), one obtains
\[
\begin{split}
U_{N,M}(t,s)P_{N}u
={}&
\sum_{a_{M}=1}^{d_{N}}
\sum_{a_{0}=1}^{d_{N}}
\mathcal{A}_{N,M}(a_{M},a_{0};t,s)
\\
&\qquad\qquad\times
\langle e_{N,a_{0}},u\rangle e_{N,a_{M}}.
\end{split}
\]
Consequently,
\[
\lim_{|\Pi_{M}|\to 0}
\mathcal{A}_{N,M}(a,b;t,s)
=
\langle e_{N,a},U_{N}(t,s)e_{N,b}\rangle
\]
for each fixed \(N\) and \(a,b\).


\subsection{Exact configuration-space kernel formula}
\label{subsec:configuration-space-kernel}

A genuine integral representation is available when the ambient Hilbert space
has a configuration-space realization. Let \((X,\mu)\) be a
\(\sigma\)-finite measure space and assume that
\[
\mathcal{H}=L^{2}(X,\mu).
\]
Suppose that the range of \(P_{N}\) is spanned by an orthonormal family
\[
\varphi_{N,1},\ldots,\varphi_{N,d_{N}}.
\]
The integral kernel of \(P_{N}\) is then
\[
\Pi_{N}(x,y)
=
\sum_{a=1}^{d_{N}}
\varphi_{N,a}(x)\overline{\varphi_{N,a}(y)}.
\]

For each time slice, the finite-rank operator
\[
E_{N,j}
=
\exp\bigl(-\mathrm{i}\Delta t_{j}H_{N}(t_{j})\bigr)
\]
has the kernel
\[
K_{N,j}(x,y)
=
\sum_{a,b=1}^{d_{N}}
E_{N,j}(a,b)
\varphi_{N,a}(x)\overline{\varphi_{N,b}(y)}.
\]

\begin{proposition}[Exact time-sliced kernel formula]
\label{prop:exact-time-sliced-kernel}
Let \(u\in\mathcal{H}_{N}\). Then
\[
\begin{split}
&(U_{N,M}(t,s)u)(x_{M})
\\
&\quad=
\int_{X^{M}}
\left(
\prod_{j=0}^{M-1}
K_{N,j}(x_{j+1},x_{j})
\right)
u(x_{0})\,
d\mu(x_{0})\cdots d\mu(x_{M-1})
\end{split}
\]
for almost every \(x_{M}\in X\).
\end{proposition}

\begin{proof}
Each \(E_{N,j}\) is a finite-rank integral operator with kernel \(K_{N,j}\).
The formula follows by repeated composition of the kernels. Since all the
operators involved have finite-dimensional range, the multiple integral is
well defined in the \(L^{2}\)-sense. Equivalently, expanding every kernel in
the orthonormal basis reduces the formula to the finite state-sum of
Definition~\ref{def:spectral-cutoff-state-sum}.
\end{proof}

\begin{remark}
\label{rem:dimension-configuration-space}
The integration variables \(x_{j}\) belong to the configuration space \(X\).
Their dimension, when \(X\) is a manifold, is \(\dim X\) and not
\(d_{N}=\dim\mathcal{H}_{N}\). The spectral cut-off changes the rank of the
kernels but does not change the underlying configuration space.
\end{remark}


\subsection{Oscillatory realizations}
\label{subsec:oscillatory-realizations}

The kernel formula of
Proposition~\ref{prop:exact-time-sliced-kernel} is always available in a
configuration-space realization, but it is not automatically a classical
oscillatory integral. Such a representation requires additional information
on the one-step kernels.

\begin{definition}[Oscillatory realization of a one-step kernel]
\label{def:oscillatory-one-step-kernel}
An oscillatory realization of \(K_{N,j}\) is a representation
\[
K_{N,j}(x,y)
=
\operatorname{Os}
\int_{\mathbb{R}^{q_{N,j}}}
\exp\bigl(\mathrm{i}\Phi_{N,j}(x,y,\theta)\bigr)
a_{N,j}(x,y,\theta)\,d\theta,
\]
where \(\Phi_{N,j}\) is a real-valued phase function and
\(a_{N,j}\) is an amplitude belonging to a specified symbol class.
The notation \(\operatorname{Os}\int\) denotes the corresponding
regularized oscillatory integral.
\end{definition}

This definition includes the case \(q_{N,j}=0\), in which the kernel is
already given by
\[
K_{N,j}(x,y)
=
a_{N,j}(x,y)
\exp\bigl(\mathrm{i}\Phi_{N,j}(x,y)\bigr).
\]
However, the existence of such a factorization is not asserted in the
abstract Hilbert-space setting.

\begin{proposition}[Composition of oscillatory one-step kernels]
\label{prop:composition-oscillatory-kernels}
Assume that every one-step kernel \(K_{N,j}\) admits an oscillatory
realization as in
Definition~\ref{def:oscillatory-one-step-kernel}, and that the amplitudes and
phases satisfy conditions allowing the corresponding oscillatory integrals
to be composed. Then the kernel of \(U_{N,M}(t,s)\) admits the representation
\[
\begin{split}
K_{N,M}(x_{M},x_{0})
={}&
\operatorname{Os}
\int_{X^{M-1}}
\int_{\mathbb{R}^{q_{N,0}+\cdots+q_{N,M-1}}}
\\
&\exp\left(
\mathrm{i}
\sum_{j=0}^{M-1}
\Phi_{N,j}(x_{j+1},x_{j},\theta_{j})
\right)
\\
&\times
\prod_{j=0}^{M-1}
a_{N,j}(x_{j+1},x_{j},\theta_{j})
\\
&\times
d\theta_{0}\cdots d\theta_{M-1}
\,d\mu(x_{1})\cdots d\mu(x_{M-1}).
\end{split}
\]
\end{proposition}

\begin{proof}
This follows by inserting the oscillatory representation of each one-step
kernel into the exact composition formula of
Proposition~\ref{prop:exact-time-sliced-kernel}. The phase of the composed
integral is the sum of the one-step phases, and its amplitude is the product
of the one-step amplitudes. The stated compatibility assumptions justify the
successive regularized integrations.
\end{proof}

\begin{definition}[Spectral cut-off oscillatory amplitude]
\label{def:spectral-cutoff-oscillatory-amplitude}
Under the hypotheses of
Proposition~\ref{prop:composition-oscillatory-kernels}, the resulting
oscillatory representation of \(U_{N,M}(t,s)\) is called a spectral cut-off
oscillatory amplitude and is denoted by
\[
I_{N,M}(t,s).
\]
This notation refers to the oscillatory representation, while
\(U_{N,M}(t,s)\) denotes the underlying time-sliced operator.
\end{definition}

Thus,
\[
I_{N,M}(t,s)=U_{N,M}(t,s)
\]
as operators whenever the oscillatory realization exists. The notation
records additional information about the representation, not a different
operator.

\begin{remark}[Absence of a canonical action]
\label{rem:no-canonical-action}
In a general abstract Hilbert space, there is no canonical configuration
space, phase function, or classical action associated with the matrix
\(H_{N}(t)\). The finite state-sum remains well defined, but it should not be
identified with a Feynman-type integral. A discrete action arises only after
choosing additional geometric data, for example a Schrödinger
representation, a phase-space quantization, or a coherent-state resolution
of the identity.
\end{remark}


\subsection{The two limiting procedures}
\label{subsec:two-limiting-procedures}

For \(u\in\mathcal{H}\), define
\[
u_{N,M}(t)
=
U_{N,M}(t,0)P_{N}u.
\]
For each fixed \(N\), finite-dimensional non-autonomous evolution theory gives
\[
\lim_{|\Pi_{M}|\to 0}
u_{N,M}(t)
=
U_{N}(t,0)P_{N}u.
\]
If the hypotheses of the spectral convergence theorem are satisfied, then
\[
\lim_{N\to\infty}
U_{N}(t,0)P_{N}u
=
U(t,0)u.
\]
Therefore
\[
U(t,0)u
=
\lim_{N\to\infty}
\lim_{|\Pi_{M}|\to 0}
U_{N,M}(t,0)P_{N}u.
\]

At the abstract level, this formula is a double limit of finite-dimensional
time-sliced operators, equivalently of the state-sums introduced in
Definition~\ref{def:spectral-cutoff-state-sum}. When the one-step kernels
possess oscillatory realizations, the same formula becomes
\[
U(t,0)u
=
\lim_{N\to\infty}
\lim_{|\Pi_{M}|\to 0}
I_{N,M}(t,0)P_{N}u.
\]

\begin{remark}
\label{rem:iterated-versus-joint-limit}
The preceding formula concerns an iterated limit: the time-slicing limit is
taken first at fixed spectral cut-off, and the spectral cut-off is then
removed. A joint limiting procedure requires estimates on the time-slicing
error that are quantitative in \(N\). Such estimates will be established
below under additional time-regularity assumptions on \(H(t)\).
\end{remark}
\section{Convergence of the spectral cut-off dynamics}
\label{sec:convergence-cutoff-dynamics}

In the previous section we introduced the finite-dimensional
oscillatory amplitudes \(I_{N,M}(t,s)\) associated with the spectral
cut-off Hamiltonians
\[
        H_N(t)=P_NH(t)P_N.
\]
For fixed \(N\), the time-slicing limit gives the finite-dimensional
propagator \(U_N(t,s)\).  The purpose of this section is to prove that,
under suitable \(H_0\)-relative assumptions, the cut-off propagators
converge strongly to the propagator of the original non-autonomous
Hamiltonian equation.

The main point is that the convergence is not proved at the level of
formal path integrals.  It is proved at the level of propagators, using
the spectral cut-off and Duhamel's formula.  The oscillatory
interpretation is then inherited from the finite-dimensional
approximations.

\subsection{The Hilbert scale associated with \(H_0\)}

Let \(H_0\) be the positive self-adjoint reference operator introduced
in Section~\ref{sec:cutoff-oscillatory-integrals}.  We assume throughout
this section that \(H_0\) has compact resolvent.  For \(s\geq0\), we set
\[
        \Hilb^s=D((1+H_0)^s),
        \qquad
        \|u\|_s=\|(1+H_0)^s u\|.
\]
Thus \(\Hilb^0=\Hilb\).  We also write
\[
        \Hilb^\infty=\bigcap_{s\geq0}\Hilb^s.
\]

The spectral projections
\[
        P_N=\mathbf 1_{[0,N]}(H_0)
\]
commute with \(H_0\).  Hence, for every \(s\geq0\),
\[
        \|P_Nu\|_s\leq \|u\|_s,
        \qquad u\in\Hilb^s.
\]
Moreover,
\[
        P_Nu\longrightarrow u
        \qquad\text{in }\Hilb^s
\]
for every \(u\in\Hilb^s\).

We shall use the following elementary compactness lemma several times.

\begin{lemma}[Uniform convergence on compact sets]
\label{lem:uniform-tail}
Let \(s\geq0\), and let \(K\subset\Hilb^s\) be compact.  Then
\[
        \sup_{u\in K}\|(I-P_N)u\|_s
        \longrightarrow0
\]
as \(N\to+\infty\).
\end{lemma}

\begin{proof}
For each fixed \(u\in\Hilb^s\), we know that
\[
        \|(I-P_N)u\|_s\longrightarrow0.
\]
Moreover, \(P_N\) is a contraction on \(\Hilb^s\).  Hence
\[
        \|(I-P_N)u\|_s\leq 2\|u\|_s.
\]
Let \(\varepsilon>0\).  Since \(K\) is compact, there exist
\(u_1,\ldots,u_m\in K\) such that for every \(u\in K\) one can find
\(u_j\) with
\[
        \|u-u_j\|_s<\varepsilon.
\]
For this \(u_j\), we have
\[
\begin{aligned}
        \|(I-P_N)u\|_s
        &\leq
        \|(I-P_N)(u-u_j)\|_s
        +
        \|(I-P_N)u_j\|_s    \\
        &\leq
        2\|u-u_j\|_s
        +
        \|(I-P_N)u_j\|_s    \\
        &<
        2\varepsilon
        +
        \|(I-P_N)u_j\|_s .
\end{aligned}
\]
Since the set \(\{u_1,\ldots,u_m\}\) is finite, there exists
\(N_0\) such that
\[
        \|(I-P_N)u_j\|_s<\varepsilon
        \qquad
        \text{for all }j=1,\ldots,m
\]
whenever \(N\geq N_0\).  Therefore
\[
        \sup_{u\in K}\|(I-P_N)u\|_s
        \leq 3\varepsilon
\]
for \(N\geq N_0\).  This proves the claim.
\end{proof}

\subsection{Assumptions on the full dynamics}

We now introduce the assumptions under which the spectral cut-off
dynamics converges.

\begin{assumption}[Well-posed Hamiltonian evolution]
\label{ass:full-dynamics}
Let \(I\subset\mathbb R\) be a compact interval.  We assume that
\(H(t)\), \(t\in I\), is a family of symmetric operators on \(\Hilb\)
such that the non-autonomous Schr\"odinger equation
\[
        \ii\partial_t u(t)=H(t)u(t)
\]
generates a unitary propagator
\[
        U(t,s):\Hilb\to\Hilb,
        \qquad s,t\in I.
\]
Thus
\[
        U(t,r)U(r,s)=U(t,s),
        \qquad
        U(s,s)=\Id,
\]
and \(U(t,s)\) is strongly continuous in \((t,s)\).

We also assume that there exists a number \(\mu\geq0\) such that:

\begin{enumerate}[label=\textup{(\roman*)}]
\item
For every \(t\in I\), the operator \(H(t)\) extends to a bounded
operator
\[
        H(t):\Hilb^\mu\longrightarrow \Hilb.
\]

\item
The map
\[
        t\longmapsto H(t)
\]
is continuous from \(I\) to
\[
        \mathcal L(\Hilb^\mu,\Hilb).
\]

\item
The propagator preserves \(\Hilb^\mu\), and for every \(u\in\Hilb^\mu\)
the map
\[
        (t,s)\longmapsto U(t,s)u
\]
is continuous from \(I\times I\) to \(\Hilb^\mu\).

\item
For every \(u\in\Hilb^\mu\), the function
\[
        t\longmapsto U(t,s)u
\]
is a strong solution of
\[
        \ii\partial_t u(t)=H(t)u(t)
\]
with values in \(\Hilb\).
\end{enumerate}
\end{assumption}

\begin{remark}
The assumption above is not meant to be optimal.  It isolates the
minimal mechanism needed for the convergence proof: the full propagator
must be defined, it must preserve enough regularity to make \(H(t)u(t)\)
meaningful, and the high spectral tail of \(U(t,s)u\) must vanish in the
\(\Hilb^\mu\)-norm.
\end{remark}

\subsection{Finite-dimensional cut-off propagators}

For each \(N\), we define
\[
        \Hilb_N=P_N\Hilb,
        \qquad
        H_N(t)=P_NH(t)P_N.
\]
Since \(\Hilb_N\) is finite-dimensional, \(H_N(t)\) is a bounded
operator on \(\Hilb_N\).

\begin{lemma}[Finite-dimensional unitary propagator]
\label{lem:finite-propagator}
For each \(N\), the cut-off equation
\[
        \ii\partial_t u_N(t)=H_N(t)u_N(t),
        \qquad
        u_N(s)=u_{N,s}\in\Hilb_N,
\]
has a unique unitary propagator
\[
        U_N(t,s):\Hilb_N\to\Hilb_N.
\]
Moreover,
\[
        U_N(t,r)U_N(r,s)=U_N(t,s),
        \qquad
        U_N(s,s)=\Id_{\Hilb_N}.
\]
\end{lemma}

\begin{proof}
Since \(H(t):\Hilb^\mu\to\Hilb\) is norm-continuous and
\(P_N\Hilb\subset \Hilb^\infty\subset \Hilb^\mu\), the map
\[
        t\longmapsto H_N(t)=P_NH(t)P_N
\]
is a continuous map from \(I\) to
\(\mathcal L(\Hilb_N,\Hilb_N)\).  Hence the finite-dimensional linear
ordinary differential equation
\[
        \partial_t u_N(t)=-\ii H_N(t)u_N(t)
\]
has a unique global solution on \(I\).

It remains to prove unitarity.  Since \(H(t)\) is symmetric, for
\(v,w\in\Hilb_N\) we have
\[
\begin{aligned}
        \langle H_N(t)v,w\rangle
        &=
        \langle P_NH(t)P_Nv,w\rangle        \\
        &=
        \langle H(t)v,w\rangle              \\
        &=
        \langle v,H(t)w\rangle              \\
        &=
        \langle v,P_NH(t)P_Nw\rangle        \\
        &=
        \langle v,H_N(t)w\rangle.
\end{aligned}
\]
Thus \(H_N(t)\) is self-adjoint on the finite-dimensional Hilbert space
\(\Hilb_N\).  If \(u_N(t)\) and \(v_N(t)\) are two solutions, then
\[
\begin{aligned}
        \frac{d}{dt}\langle u_N(t),v_N(t)\rangle
        &=
        \langle -\ii H_N(t)u_N(t),v_N(t)\rangle
        +
        \langle u_N(t),-\ii H_N(t)v_N(t)\rangle' .
\end{aligned}
\]
Using the convention that the scalar product is linear in the first
variable, the second term is
\[
        \langle u_N(t),-\ii H_N(t)v_N(t)\rangle'
        =
        \ii\langle u_N(t),H_N(t)v_N(t)\rangle.
\]
Therefore
\[
        \frac{d}{dt}\langle u_N(t),v_N(t)\rangle
        =
        -\ii\langle H_N(t)u_N(t),v_N(t)\rangle
        +
        \ii\langle u_N(t),H_N(t)v_N(t)\rangle
        =
        0.
\]
The propagator is thus unitary.  The composition law follows from
uniqueness of solutions.
\end{proof}

\begin{remark}
If the convention for the Hilbert scalar product is conjugate-linear in
the first variable, the signs in the computation above are reversed,
but the conclusion is of course the same.
\end{remark}

\subsection{The Duhamel comparison formula}

We now compare the exact propagator \(U(t,s)\) with the cut-off
propagator \(U_N(t,s)\).

Let \(u\in\Hilb^\mu\).  We define
\[
        w_N(t)=P_NU(t,s)u,
        \qquad
        v_N(t)=U_N(t,s)P_Nu.
\]
Both functions take values in \(\Hilb_N\).  The function \(v_N\)
satisfies
\[
        \ii\partial_t v_N(t)=H_N(t)v_N(t),
        \qquad
        v_N(s)=P_Nu.
\]
The projected exact solution \(w_N\) satisfies a perturbed cut-off
equation.

\begin{lemma}[Projected exact dynamics]
\label{lem:projected-exact}
For \(u\in\Hilb^\mu\), one has
\[
        \ii\partial_t w_N(t)
        =
        H_N(t)w_N(t)
        +
        R_N(t,s)u,
\]
where
\[
        R_N(t,s)u
        =
        P_NH(t)(I-P_N)U(t,s)u.
\]
\end{lemma}

\begin{proof}
Since \(U(t,s)u\) is a strong solution,
\[
        \ii\partial_t U(t,s)u=H(t)U(t,s)u.
\]
Applying \(P_N\), we get
\[
        \ii\partial_t P_NU(t,s)u
        =
        P_NH(t)U(t,s)u.
\]
Now decompose
\[
        U(t,s)u=P_NU(t,s)u+(I-P_N)U(t,s)u.
\]
Then
\[
\begin{aligned}
        P_NH(t)U(t,s)u
        &=
        P_NH(t)P_NU(t,s)u
        +
        P_NH(t)(I-P_N)U(t,s)u     \\
        &=
        H_N(t)w_N(t)
        +
        R_N(t,s)u.
\end{aligned}
\]
This proves the formula.
\end{proof}

The difference \(w_N-v_N\) is therefore controlled by the remainder
\(R_N(t,s)u\).

\begin{proposition}[Duhamel formula for the cut-off error]
\label{prop:duhamel-cutoff}
For every \(u\in\Hilb^\mu\) and every \(s,t\in I\),
\[
        P_NU(t,s)u-U_N(t,s)P_Nu
        =
        -\ii
        \int_s^t
        U_N(t,\tau)
        P_NH(\tau)(I-P_N)U(\tau,s)u
        \,d\tau .
\]
Consequently,
\[
        \|P_NU(t,s)u-U_N(t,s)P_Nu\|
        \leq
        \int_s^t
        \|H(\tau)(I-P_N)U(\tau,s)u\|
        \,d\tau
\]
when \(t\geq s\), and the analogous estimate holds with the integral
over \([t,s]\) when \(t<s\).
\end{proposition}

\begin{proof}
Set
\[
        e_N(t)=w_N(t)-v_N(t).
\]
By Lemma~\ref{lem:projected-exact} and by the equation satisfied by
\(v_N\),
\[
        \ii\partial_t e_N(t)
        =
        H_N(t)e_N(t)+R_N(t,s)u.
\]
Moreover,
\[
        e_N(s)=P_Nu-P_Nu=0.
\]
The variation-of-constants formula in the finite-dimensional space
\(\Hilb_N\) gives
\[
        e_N(t)
        =
        -\ii
        \int_s^t
        U_N(t,\tau)R_N(\tau,s)u
        \,d\tau .
\]
Substituting the definition of \(R_N\), we obtain
\[
        e_N(t)
        =
        -\ii
        \int_s^t
        U_N(t,\tau)
        P_NH(\tau)(I-P_N)U(\tau,s)u
        \,d\tau.
\]
This proves the identity.

Since \(U_N(t,\tau)\) is unitary on \(\Hilb_N\) and \(P_N\) is an
orthogonal projection on \(\Hilb\),
\[
\begin{aligned}
        \|e_N(t)\|
        &\leq
        \int_s^t
        \|U_N(t,\tau)
        P_NH(\tau)(I-P_N)U(\tau,s)u\|
        \,d\tau                                    \\
        &\leq
        \int_s^t
        \|H(\tau)(I-P_N)U(\tau,s)u\|
        \,d\tau .
\end{aligned}
\]
This gives the estimate for \(t\geq s\).  The case \(t<s\) is identical,
with the orientation of the integral reversed.
\end{proof}

\subsection{Strong convergence of the cut-off propagators}

We now prove the main theorem of this section.

\begin{theorem}[Strong convergence of the spectral cut-off dynamics]
\label{thm:strong-cutoff-convergence}
Assume Assumption~\ref{ass:full-dynamics}.  Then, for every
\(u\in\Hilb^\mu\),
\[
        U_N(t,s)P_Nu
        \longrightarrow
        U(t,s)u
\]
strongly in \(\Hilb\), uniformly for \((t,s)\in I\times I\).
\end{theorem}

\begin{proof}
We write
\[
\begin{aligned}
        U_N(t,s)P_Nu-U(t,s)u
        &=
        \bigl(U_N(t,s)P_Nu-P_NU(t,s)u\bigr)
        +
        \bigl(P_NU(t,s)u-U(t,s)u\bigr).
\end{aligned}
\]
We estimate the two terms separately.

First consider the projection tail
\[
        (P_N-I)U(t,s)u.
\]
By Assumption~\ref{ass:full-dynamics}, the map
\[
        (t,s)\longmapsto U(t,s)u
\]
is continuous from the compact set \(I\times I\) to \(\Hilb^\mu\).
Hence
\[
        K_u=\{U(t,s)u;(t,s)\in I\times I\}
\]
is compact in \(\Hilb^\mu\).  By Lemma~\ref{lem:uniform-tail},
\[
        \sup_{(t,s)\in I\times I}
        \|(I-P_N)U(t,s)u\|_\mu
        \longrightarrow0.
\]
In particular,
\[
        \sup_{(t,s)\in I\times I}
        \|(I-P_N)U(t,s)u\|
        \longrightarrow0.
\]

We now estimate the first term.  By
Proposition~\ref{prop:duhamel-cutoff},
\[
\begin{aligned}
        \|P_NU(t,s)u-U_N(t,s)P_Nu\|
        &\leq
        \int_{\min(s,t)}^{\max(s,t)}
        \|H(\tau)(I-P_N)U(\tau,s)u\|
        \,d\tau .
\end{aligned}
\]
Since \(t\mapsto H(t)\) is continuous from \(I\) to
\(\mathcal L(\Hilb^\mu,\Hilb)\), there exists a constant \(C_I>0\) such
that
\[
        \|H(\tau)z\|
        \leq C_I\|z\|_\mu
\]
for all \(\tau\in I\) and all \(z\in\Hilb^\mu\).  Therefore
\[
\begin{aligned}
        \|P_NU(t,s)u-U_N(t,s)P_Nu\|
        &\leq
        C_I |t-s|
        \sup_{\tau\in I}
        \|(I-P_N)U(\tau,s)u\|_\mu .
\end{aligned}
\]
Again, the set
\[
        \{U(\tau,s)u;(\tau,s)\in I\times I\}
\]
is compact in \(\Hilb^\mu\).  Hence, by
Lemma~\ref{lem:uniform-tail},
\[
        \sup_{(\tau,s)\in I\times I}
        \|(I-P_N)U(\tau,s)u\|_\mu
        \longrightarrow0.
\]
It follows that
\[
        \sup_{(t,s)\in I\times I}
        \|P_NU(t,s)u-U_N(t,s)P_Nu\|
        \longrightarrow0.
\]
Combining this with the projection-tail estimate proves
\[
        \sup_{(t,s)\in I\times I}
        \|U_N(t,s)P_Nu-U(t,s)u\|
        \longrightarrow0.
\]
This is the desired uniform strong convergence.
\end{proof}

\begin{remark}
The proof shows precisely where the spectral cut-off enters.  The only
error term is
\[
        P_NH(t)(I-P_N)U(t,s)u.
\]
Thus the convergence follows from the fact that the high-energy tail of
the exact solution is small in the regularity norm required to apply
\(H(t)\).
\end{remark}

\subsection{Consequences for state-sum and oscillatory amplitudes}

Recall that for fixed \(N\), the time-sliced operator \(U_{N,M}(t,s)\)
converges to \(U_N(t,s)\) as the mesh of the partition tends to zero.
Combining this finite-dimensional convergence with
Theorem~\ref{thm:strong-cutoff-convergence}, we obtain the following
consequence.

\begin{corollary}[State-sum construction of the solution]
\label{cor:state-sum-construction}
Let Assumption~\ref{ass:full-dynamics} hold.  Let \(u\in\Hilb^\mu\).
For each \(N\) and each time partition, consider the time-sliced operator
\[
        U_{N,M}(t,s)
        =
        \prod_{j=M-1}^{0}
        \exp\bigl(-\ii\Delta t_jH_N(t_j)\bigr).
\]
Then
\[
        U(t,s)u
        =
        \lim_{N\to\infty}
        \lim_{M\to\infty}
        U_{N,M}(t,s)P_Nu
\]
strongly in \(\Hilb\), uniformly for \((t,s)\in I\times I\).
Equivalently, the finite-dimensional state-sum representation of
\(U_{N,M}(t,s)P_Nu\) converges to \(U(t,s)u\) in this iterated sense.
\end{corollary}

\begin{proof}
For fixed \(N\), the finite-dimensional time-slicing construction gives
\[
        \lim_{M\to\infty}
        U_{N,M}(t,s)P_Nu
        =
        U_N(t,s)P_Nu
\]
in \(\Hilb_N\).  By Theorem~\ref{thm:strong-cutoff-convergence},
\[
        \lim_{N\to\infty}
        U_N(t,s)P_Nu
        =
        U(t,s)u
\]
strongly in \(\Hilb\), uniformly for \((t,s)\in I\times I\).  Combining
the two limits gives the result.
\end{proof}

\begin{corollary}[Conditional oscillatory construction]
\label{cor:oscillatory-construction}
Under the assumptions of
Corollary~\ref{cor:state-sum-construction}, suppose in addition that the
one-step kernels admit oscillatory realizations in the sense of
Definition~\ref{def:oscillatory-one-step-kernel}. Then
\[
 U(t,s)u
 =
 \lim_{N\to\infty}\lim_{M\to\infty}
 I_{N,M}(t,s)P_Nu
\]
strongly in \(\Hilb\), uniformly for \((t,s)\in I\times I\).
\end{corollary}

\begin{proof}
By definition of an oscillatory realization,
\[
 I_{N,M}(t,s)=U_{N,M}(t,s)
\]
as operators. The conclusion therefore follows from
Corollary~\ref{cor:state-sum-construction}.
\end{proof}

\begin{remark}
The abstract construction is intrinsically a limit of finite-dimensional
state-sums. It becomes a limit of oscillatory integrals only in realizations
for which the required phases and amplitudes have been constructed.
\end{remark}

\subsection{Quantitative spectral convergence}
\label{subsec:quantitative-spectral-convergence}

The preceding convergence theorem is qualitative. We now show that additional
regularity of the exact evolution gives an explicit convergence rate with
respect to the spectral cut-off. Set
\[
\Lambda=1+H_{0}.
\]
Since
\[
P_{N}=\mathbf{1}_{[0,N]}(H_{0})
      =\mathbf{1}_{[1,1+N]}(\Lambda),
\]
the spectral theorem gives the following elementary tail estimate.

\begin{lemma}[Quantitative spectral tail]
\label{lem:quantitative-spectral-tail}
Let \(r\geq 0\) and \(\sigma\geq 0\). Then, for every
\(u\in\Hilb^{r+\sigma}\),
\[
\|(I-P_{N})u\|_{r}
\leq
(1+N)^{-\sigma}\|u\|_{r+\sigma}.
\]
\end{lemma}

\begin{proof}
By the spectral theorem,
\[
\|(I-P_{N})u\|_{r}^{2}
=
\int_{(N,\infty)}
(1+\lambda)^{2r}\,d\langle E_{H_{0}}(\lambda)u,u\rangle.
\]
For \(\lambda>N\), one has
\[
(1+\lambda)^{2r}
\leq
(1+N)^{-2\sigma}(1+\lambda)^{2(r+\sigma)}.
\]
Consequently,
\[
\|(I-P_{N})u\|_{r}^{2}
\leq
(1+N)^{-2\sigma}\|u\|_{r+\sigma}^{2},
\]
which proves the claim.
\end{proof}

We can now quantify the convergence established in
Theorem~\ref{thm:strong-cutoff-convergence}.

\begin{theorem}[Quantitative convergence of the cut-off dynamics]
\label{thm:quantitative-cutoff-convergence}
Let \(I\subset\mathbb{R}\) be compact, let \(\mu\geq 0\), and assume that
\[
H(t)\in\mathcal{L}(\Hilb^{\mu},\Hilb)
\]
depends continuously on \(t\in I\). Suppose that the non-autonomous equation
generates a unitary propagator \(U(t,s)\) and that, for some \(\sigma>0\),
\[
U(t,s)\Hilb^{\mu+\sigma}
\subset
\Hilb^{\mu+\sigma}.
\]
Assume moreover that
\[
C_{U,\mu+\sigma}
:=
\sup_{s,t\in I}
\|U(t,s)\|_{\mathcal{L}(\Hilb^{\mu+\sigma})}
<\infty.
\]
Set
\[
C_{H,\mu}
:=
\sup_{t\in I}
\|H(t)\|_{\mathcal{L}(\Hilb^{\mu},\Hilb)}.
\]
Then, for every \(u\in\Hilb^{\mu+\sigma}\),
\[
\begin{split}
\sup_{s,t\in I}
\|U_{N}(t,s)P_{N}u-U(t,s)u\|
\leq{}&
C_{U,\mu+\sigma}
\Bigl(
(1+N)^{-(\mu+\sigma)}
\\
&\quad
+
|I|C_{H,\mu}(1+N)^{-\sigma}
\Bigr)
\|u\|_{\mu+\sigma}.
\end{split}
\]
In particular, there exists a constant \(C_{I,\mu,\sigma}>0\),
independent of \(N\), such that
\[
\sup_{s,t\in I}
\|U_{N}(t,s)P_{N}u-U(t,s)u\|
\leq
C_{I,\mu,\sigma}(1+N)^{-\sigma}
\|u\|_{\mu+\sigma}.
\]
\end{theorem}

\begin{proof}
As in the proof of Theorem~\ref{thm:strong-cutoff-convergence}, write
\[
\begin{split}
U_{N}(t,s)P_{N}u-U(t,s)u
={}&
U_{N}(t,s)P_{N}u-P_{N}U(t,s)u
\\
&+
(P_{N}-I)U(t,s)u.
\end{split}
\]
The second term is estimated using
Lemma~\ref{lem:quantitative-spectral-tail} with \(r=0\):
\[
\begin{split}
\|(I-P_{N})U(t,s)u\|
&\leq
(1+N)^{-(\mu+\sigma)}
\|U(t,s)u\|_{\mu+\sigma}
\\
&\leq
C_{U,\mu+\sigma}
(1+N)^{-(\mu+\sigma)}
\|u\|_{\mu+\sigma}.
\end{split}
\]

For the first term, Proposition~\ref{prop:duhamel-cutoff} gives
\[
\begin{split}
&\|P_{N}U(t,s)u-U_{N}(t,s)P_{N}u\|
\\
&\qquad\leq
\int_{\min(s,t)}^{\max(s,t)}
\|H(\tau)(I-P_{N})U(\tau,s)u\|\,d\tau
\\
&\qquad\leq
C_{H,\mu}|t-s|
\sup_{\tau\in I}
\|(I-P_{N})U(\tau,s)u\|_{\mu}.
\end{split}
\]
Applying Lemma~\ref{lem:quantitative-spectral-tail} with \(r=\mu\)
yields
\[
\|(I-P_{N})U(\tau,s)u\|_{\mu}
\leq
(1+N)^{-\sigma}
\|U(\tau,s)u\|_{\mu+\sigma}.
\]
Therefore,
\[
\begin{split}
&\|P_{N}U(t,s)u-U_{N}(t,s)P_{N}u\|
\\
&\qquad\leq
|I|C_{H,\mu}C_{U,\mu+\sigma}
(1+N)^{-\sigma}
\|u\|_{\mu+\sigma}.
\end{split}
\]
Combining the two estimates proves the result.
\end{proof}

\begin{remark}
\label{rem:spectral-rate-interpretation}
The convergence rate is determined by the regularity of the exact orbit
relative to the reference operator \(H_{0}\). In concrete elliptic settings,
the parameter \(\sigma\) measures the number of additional Sobolev derivatives
of the initial datum and of the propagated solution. The estimate is therefore
the natural spectral analogue of the usual rate of Galerkin approximation.
\end{remark}


\subsection{Uniform stability from commutator estimates}
\label{subsec:commutator-stability}

The quantitative theorem above still uses regularity bounds for the exact
propagator. We next give an intrinsic sufficient condition for the uniform
stability of the cut-off propagators. This condition is formulated directly
in terms of the commutator of \(H(t)\) with the Hilbert scale generated by
\(\Lambda=1+H_{0}\).

\begin{assumption}[Energy commutator bound]
\label{ass:energy-commutator-bound}
Let \(r\geq 0\). Assume that there exists a locally integrable function
\(c_{r}:I\to[0,\infty)\) such that
\[
\left|
\ii
\left(
\langle H(t)v,\Lambda^{2r}v\rangle
-
\langle \Lambda^{2r}v,H(t)v\rangle
\right)
\right|
\leq
c_{r}(t)\|v\|_{r}^{2}
\]
for every \(v\in\Hilb^{\infty}\) and almost every \(t\in I\).
\end{assumption}

Equivalently, whenever the commutator is defined as a quadratic form,
Assumption~\ref{ass:energy-commutator-bound} reads
\[
\left|
\langle v,\ii[H(t),\Lambda^{2r}]v\rangle
\right|
\leq
c_{r}(t)\|v\|_{r}^{2}.
\]

\begin{proposition}[Uniform stability of the cut-off propagators]
\label{prop:uniform-cutoff-stability}
Assume Assumption~\ref{ass:energy-commutator-bound}. Then, for every \(N\),
every \(v\in\Hilb_{N}\), and every \(s,t\in I\),
\[
\|U_{N}(t,s)v\|_{r}
\leq
\exp\left(
\frac{1}{2}
\int_{\min(s,t)}^{\max(s,t)}c_{r}(\tau)\,d\tau
\right)
\|v\|_{r}.
\]
In particular, the stability constant is independent of \(N\).
\end{proposition}

\begin{proof}
Fix \(N\), \(s\in I\), and \(v\in\Hilb_{N}\), and set
\[
v_{N}(t)=U_{N}(t,s)v.
\]
Since \(\Hilb_{N}\subset\Hilb^{\infty}\), all the following
computations take place in a finite-dimensional space. We have
\[
\ii\partial_{t}v_{N}(t)=H_{N}(t)v_{N}(t).
\]
Hence
\[
\begin{split}
\frac{d}{dt}\|v_{N}(t)\|_{r}^{2}
&=
\frac{d}{dt}
\langle \Lambda^{r}v_{N}(t),\Lambda^{r}v_{N}(t)\rangle
\\
&=
\ii
\left(
\langle H_{N}(t)v_{N}(t),\Lambda^{2r}v_{N}(t)\rangle
\right.
\\
&\hspace{35mm}\left.
-
\langle \Lambda^{2r}v_{N}(t),H_{N}(t)v_{N}(t)\rangle
\right).
\end{split}
\]
Because \(P_{N}\) commutes with \(\Lambda\) and
\(v_{N}(t)\in\Hilb_{N}\), the expression on the right-hand side equals
\[
\ii
\left(
\langle H(t)v_{N}(t),\Lambda^{2r}v_{N}(t)\rangle
-
\langle \Lambda^{2r}v_{N}(t),H(t)v_{N}(t)\rangle
\right).
\]
Assumption~\ref{ass:energy-commutator-bound} therefore implies
\[
\left|
\frac{d}{dt}\|v_{N}(t)\|_{r}^{2}
\right|
\leq
c_{r}(t)\|v_{N}(t)\|_{r}^{2}.
\]
Gronwall's inequality gives
\[
\|v_{N}(t)\|_{r}^{2}
\leq
\exp\left(
\int_{\min(s,t)}^{\max(s,t)}c_{r}(\tau)\,d\tau
\right)
\|v\|_{r}^{2}.
\]
Taking square roots proves the result.
\end{proof}

\begin{remark}
The important point is that the estimate is obtained before passing to the
limit and is uniform in the spectral cut-off. Thus the stability assumption
used in the energy-norm convergence theorem is no longer merely postulated:
it follows from a commutator estimate on the original Hamiltonian family.
\end{remark}


\subsection{Construction of the propagator by removal of the cut-off}
\label{subsec:constructive-propagator}

We now remove the assumption that the full propagator is known in advance.
Under uniform estimates on the spectral scale, the cut-off propagators form a
Cauchy family and therefore construct the non-autonomous evolution.


\begin{theorem}[Construction by spectral cut-off]
\label{thm:construction-by-spectral-cutoff}
Let \(I\subset\mathbb{R}\) be compact. Assume that there exist
\(\mu\geq 0\) and \(\sigma>0\) such that the following conditions hold.

\begin{enumerate}
\item For every \(r\geq 0\), the family \(H(t)\) is norm-continuous as a map
\[
I\longrightarrow
\mathcal{L}(\Hilb^{r+\mu},\Hilb^{r}).
\]

\item The energy commutator bound of
Assumption~\ref{ass:energy-commutator-bound} holds for
\(r=\mu+\sigma\), with
\[
c_{\mu+\sigma}\in L^{1}(I).
\]

\item Each \(H(t)\) is symmetric on the common core
\(\Hilb^{\infty}\).
\end{enumerate}

Then, for every \(u\in\Hilb^{\mu+\sigma}\), the family
\[
U_{N}(t,s)P_{N}u
\]
is Cauchy in \(\Hilb\), uniformly for \((t,s)\in I\times I\).
Consequently, the limit
\[
U(t,s)u
:=
\lim_{N\to\infty}U_{N}(t,s)P_{N}u
\]
exists uniformly on \(I\times I\).

Moreover, there exists a constant \(C_{I,\mu,\sigma}>0\) such that, whenever
\(M\geq N\),
\[
\sup_{s,t\in I}
\|U_{M}(t,s)P_{M}u-U_{N}(t,s)P_{N}u\|
\leq
C_{I,\mu,\sigma}(1+N)^{-\sigma}
\|u\|_{\mu+\sigma}.
\]
The limiting family extends uniquely to a strongly continuous unitary
propagator on \(\Hilb\).

If the preceding assumptions hold at every level of the Hilbert scale, then
\(U(t,s)\) preserves \(\Hilb^{\infty}\), and for
\(u\in\Hilb^{\infty}\) one has
\[
\ii\partial_{t}U(t,s)u
=
H(t)U(t,s)u.
\]
Thus the full non-autonomous Hamiltonian evolution is constructed as the
spectral cut-off limit.
\end{theorem}

\begin{proof}
Fix \(M\geq N\), \(s\in I\), and
\(u\in\Hilb^{\mu+\sigma}\). Set
\[
w_{M}(t)=U_{M}(t,s)P_{M}u
\]
and
\[
z_{N,M}(t)=P_{N}w_{M}(t).
\]
Since \(P_{N}P_{M}=P_{N}\), we have
\[
z_{N,M}(s)=P_{N}u.
\]
Moreover,
\[
\begin{split}
\ii\partial_{t}z_{N,M}(t)
&=
P_{N}H(t)P_{M}w_{M}(t)
\\
&=
P_{N}H(t)P_{N}z_{N,M}(t)
+
P_{N}H(t)(P_{M}-P_{N})w_{M}(t)
\\
&=
H_{N}(t)z_{N,M}(t)
+
P_{N}H(t)(P_{M}-P_{N})w_{M}(t).
\end{split}
\]
Variation of constants in \(\Hilb_{N}\) gives
\[
\begin{split}
z_{N,M}(t)-U_{N}(t,s)P_{N}u
=
-\ii
\int_{s}^{t}
U_{N}(t,\tau)
P_{N}H(\tau)(P_{M}-P_{N})w_{M}(\tau)\,d\tau.
\end{split}
\]
By unitarity of \(U_{N}\) in \(\Hilb\),
\[
\begin{split}
&\|z_{N,M}(t)-U_{N}(t,s)P_{N}u\|
\\
&\qquad\leq
C_{H,\mu}|I|
\sup_{\tau\in I}
\|(P_{M}-P_{N})w_{M}(\tau)\|_{\mu},
\end{split}
\]
where
\[
C_{H,\mu}
=
\sup_{\tau\in I}
\|H(\tau)\|_{\mathcal{L}(\Hilb^{\mu},\Hilb)}.
\]
Since
\[
P_{M}-P_{N}=P_{M}(I-P_{N}),
\]
Lemma~\ref{lem:quantitative-spectral-tail} implies
\[
\|(P_{M}-P_{N})w_{M}(\tau)\|_{\mu}
\leq
(1+N)^{-\sigma}\|w_{M}(\tau)\|_{\mu+\sigma}.
\]
By Proposition~\ref{prop:uniform-cutoff-stability},
\[
\sup_{\tau,s\in I}
\|U_{M}(\tau,s)P_{M}u\|_{\mu+\sigma}
\leq
C_{\mathrm{stab}}\|u\|_{\mu+\sigma},
\]
where
\[
C_{\mathrm{stab}}
=
\exp\left(
\frac{1}{2}\int_{I}c_{\mu+\sigma}(\tau)\,d\tau
\right)
\]
is independent of \(M\). Hence
\[
\sup_{s,t\in I}
\|z_{N,M}(t)-U_{N}(t,s)P_{N}u\|
\leq
C_{H,\mu}|I|C_{\mathrm{stab}}
(1+N)^{-\sigma}
\|u\|_{\mu+\sigma}.
\]

It remains to estimate the component of \(w_{M}(t)\) orthogonal to
\(\Hilb_{N}\). By Lemma~\ref{lem:quantitative-spectral-tail},
\[
\begin{split}
\|(I-P_{N})w_{M}(t)\|
&\leq
(1+N)^{-(\mu+\sigma)}
\|w_{M}(t)\|_{\mu+\sigma}
\\
&\leq
C_{\mathrm{stab}}
(1+N)^{-(\mu+\sigma)}
\|u\|_{\mu+\sigma}.
\end{split}
\]
Since
\[
w_{M}(t)-U_{N}(t,s)P_{N}u
=
(I-P_{N})w_{M}(t)
+
z_{N,M}(t)-U_{N}(t,s)P_{N}u,
\]
we obtain the asserted Cauchy estimate.

The limit therefore exists for
\(u\in\Hilb^{\mu+\sigma}\). Since this space is dense in
\(\Hilb\) and
\[
\|U_{N}(t,s)P_{N}u\|=\|P_{N}u\|\leq\|u\|,
\]
the limiting operators extend uniquely to contractions on \(\Hilb\).
Applying the same construction to the backward propagators \(U_{N}(s,t)\)
shows that
\[
U(t,s)U(s,t)=U(s,t)U(t,s)=I.
\]
Thus \(U(t,s)\) is unitary.

The propagator identity follows by passing to the limit in
\[
U_{N}(t,r)U_{N}(r,s)=U_{N}(t,s),
\]
first on the dense regularity space and then on \(\Hilb\).
Strong continuity follows from the uniform convergence on \(I\times I\).

Finally, if the assumptions hold at every level of the Hilbert scale, the
same argument gives convergence in each \(\Hilb^{r}\). For
\(u\in\Hilb^{\infty}\), one may therefore pass to the limit in the
integral equation
\[
U_{N}(t,s)P_{N}u
=
P_{N}u
-
\ii
\int_{s}^{t}
P_{N}H(\tau)P_{N}
U_{N}(\tau,s)P_{N}u\,d\tau.
\]
This yields
\[
U(t,s)u
=
u
-
\ii
\int_{s}^{t}H(\tau)U(\tau,s)u\,d\tau.
\]
The norm-continuity of
\[
t\longmapsto H(t)\in
\mathcal{L}(\Hilb^{r+\mu},\Hilb^{r})
\]
then implies that \(t\mapsto U(t,s)u\) is differentiable in
\(\Hilb\) and satisfies
\[
\ii\partial_{t}U(t,s)u=H(t)U(t,s)u.
\]
\end{proof}

\begin{corollary}[Quantitative convergence in energy norms]
\label{cor:quantitative-energy-convergence}
Let \(r\geq0\) and \(\sigma>0\). Assume that
\[
H(t):\Hilb^{r+\mu}\longrightarrow\Hilb^r
\]
is norm-continuous in \(t\), and that the energy commutator bound holds
at both levels \(r\) and \(r+\mu+\sigma\). Then there exists
\(C_{I,r,\mu,\sigma}>0\) such that
\[
\sup_{s,t\in I}
\|U_N(t,s)P_Nu-U(t,s)u\|_r
\leq
C_{I,r,\mu,\sigma}(1+N)^{-\sigma}
\|u\|_{r+\mu+\sigma}
\]
for every \(u\in\Hilb^{r+\mu+\sigma}\).
\end{corollary}

\end{document}
